\documentclass[11pt,a4paper,french]{imta}

\title{Rapport : Projet robot aspirateur}

\author{
\tDAHANI Wissal \\
\tLOUMOUAMOU Yannis\\
\tMAHMOUDI Zaineb\\
}

\day17
\month03
\year2026

% Mode de compilation robuste (VM / environnements incomplets)
\setkeys{Gin}{draft}
\makeatletter
\renewcommand{\IMTAfrontcover}{}
\renewcommand{\IMTAcoverpage}{}
\makeatother

\begin{document}

\section{Introduction}
Ce projet vise a realiser un robot aspirateur autonome en simulation. Le sujet pedagogique de reference est fonde sur des chapitres ROS1 (7 a 10) de \emph{Programming Robots with ROS}. Notre equipe a choisi de relever le defi d'une implementation complete en ROS2 Humble.

Objectifs fonctionnels:
\begin{enumerate}
  \item Evitement reactif d'obstacles
  \item Cartographie SLAM d'un environnement inconnu
  \item Navigation sur carte avec Nav2
  \item Couverture de zone de type aspirateur
\end{enumerate}

\section{Contexte de migration ROS1 vers ROS2}
La migration ROS1 vers ROS2 a implique les changements suivants:
\begin{itemize}
  \item \texttt{rospy} vers \texttt{rclpy}
  \item \texttt{catkin} vers \texttt{ament/colcon}
  \item \texttt{move\_base} vers Nav2
  \item \texttt{gmapping} vers \texttt{slam\_toolbox}
\end{itemize}

Cette migration constitue un apport technique majeur du projet: elle preserve la logique des chapitres du livre tout en s'appuyant sur une pile logicielle recente.

\section{Environnement experimental}
\subsection{Plateforme}
\begin{itemize}
  \item Hote: Windows
  \item Virtualisation: VirtualBox
  \item Systeme invite: Ubuntu 22.04
  \item ROS: ROS2 Humble
  \item Simulation: Gazebo + TurtleBot3 Burger
\end{itemize}

\subsection{Paquets utilises}
\begin{itemize}
  \item Paquets developpes: \texttt{reactive\_avoidance}, \texttt{coverage\_planner}
  \item Paquets externes: \texttt{nav2\_bringup}, \texttt{navigation2}, \texttt{slam\_toolbox}
\end{itemize}

\section{Architecture mise en oeuvre}
\subsection{Controle reactif (chapitres 7-8)}
Le noeud \texttt{reactive\_avoidance} lit \texttt{/scan} (LaserScan) et publie \texttt{/cmd\_vel} (Twist).

Ameliorations apportees pour stabiliser le comportement:
\begin{itemize}
  \item Hysteresis sur le seuil de decision obstacle
  \item Boucle de commande a frequence fixe
  \item Zone prudente + zone d'urgence pour anticiper les collisions
  \item Reduction des vitesses par defaut en VM
\end{itemize}

\subsection{Cartographie et navigation (chapitres 9-10)}
\begin{itemize}
  \item Cartographie: \texttt{slam\_toolbox} avec sauvegarde \texttt{.yaml/.pgm}
  \item Navigation: \texttt{nav2\_bringup} avec configuration dediee
  \item Localisation: AMCL + initialisation de pose dans RViz2
\end{itemize}

\subsection{Couverture de zone}
Le noeud \texttt{coverage\_planner} genere des objectifs de couverture (pattern lawnmower) et les publie sequentiellement sur \texttt{/goal\_pose} en suivant la progression via \texttt{/amcl\_pose}.

\section{Etat d'avancement actuel}
\begin{center}
\begin{tabular}{p{3.6cm}p{2.8cm}p{7.2cm}}
\toprule
\textbf{Bloc} & \textbf{Etat} & \textbf{Commentaires} \\
\midrule
Chapitres 7-8 & En cours avancee & Publication \texttt{/cmd\_vel} confirmee; stabilite physique en cours d'optimisation en simulation \\
Chapitre 9 & Pret a valider & Configuration SLAM en place; validation finale via carte sauvegardee \\
Chapitre 10 & En cours & Nav2 demarre; ajustements TF/AMCL et procedure d'initial pose en cours \\
Couverture & Prototype fonctionnel & Generation et publication sequentielle d'objectifs implementees \\
\bottomrule
\end{tabular}
\end{center}

\section{Obstacles rencontres et resolutions}
\subsection{Packages non trouves}
\textbf{Symptome:} \texttt{Package not found} (\texttt{nav2\_bringup}, \texttt{coverage\_planner}).

\textbf{Resolution:} installation des dependances, build \texttt{colcon}, puis double \texttt{source}:
\begin{bash}
source /opt/ros/humble/setup.bash
source ~/ros2_ws/install/setup.bash
\end{bash}

\subsection{Parametrage Nav2 incomplet}
\textbf{Symptome:} erreurs \texttt{yaml\_filename not initialized} et \texttt{No critics defined for FollowPath}.

\textbf{Resolution:} fichier \texttt{nav2\_params.yaml} complete (plugins, critics DWB, frames coherentes).

\subsection{Erreurs TF (map/odom/base\_link)}
\textbf{Symptome:} \texttt{Invalid frame ID} et timeouts costmaps.

\textbf{Resolution:}
\begin{itemize}
  \item Harmonisation sur \texttt{base\_footprint}
  \item Verification des transformations avec \texttt{tf2\_echo}
  \item Initialisation explicite de pose AMCL dans RViz2
\end{itemize}

\subsection{Instabilite physique en simulation}
\textbf{Symptome:} le robot percute, chute, puis ne se releve pas.

\textbf{Resolution:}
\begin{itemize}
  \item Vitesse lineaire et angulaire reduites
  \item Distances de securite augmentees
  \item Controle reactif adouci (hysteresis + zones prudente/urgence)
\end{itemize}

\section{Plan de validation final}
\begin{enumerate}
  \item Valider completement chapitres 7-8 (preuve avance + preuve rotation)
  \item Valider chapitre 9 (carte exploitable sauvegardee)
  \item Valider chapitre 10 (3 objectifs atteints sans collision)
  \item Evaluer la couverture avec metriques quantitatives
\end{enumerate}

\section{Conclusion}
Le projet a deja atteint une base technique solide sous ROS2, malgre un sujet initialement pense pour ROS1. Les principaux risques restants sont maintenant concentres sur la robustesse experimentale (physique simulation, TF et tuning Nav2). La suite du travail porte surtout sur la validation chiffrable et la consolidation des resultats pour la version finale du rapport.

\newpage
\begin{IMTAannexes}

\IMTAannexe{Commandes de reference}

\section{Build et environnement}
\begin{bash}
source /opt/ros/humble/setup.bash
cd ~/ros2_ws
colcon build --symlink-install --base-paths src/ROS2Project_ROBIN_2026/ros2_pkgs
source ~/ros2_ws/install/setup.bash
\end{bash}

\section{Validation chapitres 7-8}
\begin{bash}
export TURTLEBOT3_MODEL=burger
ros2 launch turtlebot3_gazebo turtlebot3_world.launch.py
ros2 launch reactive_avoidance avoidance.launch.py
ros2 topic echo /cmd_vel
\end{bash}

\section{SLAM et navigation}
\begin{bash}
export PROJECT_ROOT=$HOME/ros2_ws/src/ROS2Project_ROBIN_2026
ros2 launch slam_toolbox online_async_launch.py \
  slam_params_file:=$PROJECT_ROOT/config/slam_params.yaml \
  use_sim_time:=true

ros2 launch nav2_bringup bringup_launch.py \
  use_sim_time:=true \
  map:=$HOME/ros2_ws/maps/my_map.yaml \
  params_file:=$PROJECT_ROOT/config/nav2_params.yaml
\end{bash}

\IMTAannexe{Journal de mise a jour}
Ce rapport est evolutif et sera mis a jour apres chaque campagne de test.

\begin{itemize}
  \item Version actuelle: architecture stabilisee, validation experimentale en cours
  \item Prochaine mise a jour: resultats quantitatifs (couverture, temps, collisions)
\end{itemize}

\end{IMTAannexes}

\end{document}
